\documentclass[11pt]{article}

\usepackage[margin=1in]{geometry}
\usepackage{amsmath,amssymb,amsthm,mathtools}
\usepackage{enumitem}
\usepackage{microtype}
\usepackage[hidelinks]{hyperref}
\usepackage{xurl}

\newtheorem{theorem}{Theorem}
\newtheorem{remark}{Remark}

\newcommand{\MaxL}{\mathcal{L}}

\title{A Tiered Deployment Pattern for Reader Privacy:\\
\large Channel-Bound Primary and Fallback Paths}
\author{}
\date{June 18, 2026 (archive rev0900)}

\begin{document}
\maketitle

\begin{abstract}
A primary cryptographic path plus a residual DHT fallback can be a useful deployment pattern, but the previous version overstated what scalar observation and fallback rates establish. Dummy-width pressure is exact only inside a proved channel model, and the overall two-path budget is valid only when each path has a source-bound channel bound and the path selector is either secret-independent or explicitly included as a witness. This revision keeps the architecture while replacing scalar \(p_{\mathrm{obs}}\) receipts with channel-witness records. A one-bit counterexample shows why: even if both path outputs are individually uninformative, choosing the path as \(W=C\) reveals the secret perfectly. The fallback is therefore ``honest'' only when path selection, observation, retries, stopping, and drift are all bound to the receipt.
\end{abstract}

\paragraph{Scope.}
This note exports an architecture, not a deployment anonymity certificate. Boss Fight owns repeated-use and control-model arithmetic; the observation and tier notes own channel wrappers and joint composition; the line-item note owns the channel-bound receipt carrier.\cite{series:bossfightB,series:obsatten,series:tieredobs,series:receiptitems}

\section{Why retain a two-path architecture}\label{sec:why}
A managed primary surface can make expensive query-privacy mechanisms, auditing, and change control tractable. A DHT fallback can preserve availability for long-tail requests and outages. The split is operationally useful even though neither path receives privacy merely from its label.

The dummy-width calculations in Boss Fight~B show that, under an explicit independent-erasure control model, tiny per-lookup allocations can force cover near the full candidate universe. Those values motivate a mechanism change; they do not establish the observation channel of IPFS, a delegated router, or a fallback implementation.\cite{series:bossfightB}

The architecture therefore pursues three separate levers:
\begin{enumerate}[leftmargin=*]
\item reduce repeated releases with caching, rate limits, or coarser queries;
\item change and prove the observer channel, for example through cryptographic query privacy or a source-bound common-mixture wrapper;
\item constrain the residual fallback and account for its selector as part of the complete transcript.
\end{enumerate}

\section{The two paths}\label{sec:paths}
\paragraph{Primary path.}
The primary path uses a small, managed interface with a declared reader-privacy mechanism. Its receipt binds the contacted service surface, query and response projection, failure/fallback semantics, source/configuration, channel witness, budget, and drift invalidators. Destination equalization, when used, is a separate channel theorem rather than a synonym for query privacy.\cite{series:mceq,series:deployedsurfaces}

\paragraph{Fallback path.}
The fallback path declares the raw-DHT projection, horizon, contact/retry/stop behavior, observer, source/configuration, and a direct or theorem-backed channel bound. A compromise fraction, contact count, or reveal rate is telemetry unless the channel witness proves how it enters the bound. Rate limiting \(Q\) is a policy input; it does not turn an unproved per-lookup value into a valid window budget.

\section{Path selection is itself a channel}\label{sec:selector}
Let \(W\in\{0,1\}\) select the primary or fallback path and let \(Y_W\) be the corresponding path transcript. If path choice or its consequences are observable, the complete output includes \(W\).

\begin{theorem}[Secret-independent path mixture]\label{thm:mixture}
Suppose \(W\) is independent of \(C\), \(P(W=1)=\pi\), and the conditional path channels satisfy
\[
\MaxL(C\to Y_w\mid W=w)\le b_w.
\]
Then the visible-path output \((W,Y_W)\) satisfies
\[
\MaxL(C\to(W,Y_W))
\le\log_2\!\left((1-\pi)2^{b_0}+\pi2^{b_1}\right).
\]
\end{theorem}
\begin{proof}
The outputs for \(W=0\) and \(W=1\) are disjoint. Their column-max sums are weighted by \(1-\pi\) and \(\pi\), respectively, and each is bounded by \(2^{b_w}\).
\end{proof}

The theorem requires a channel-level independence statement. A monitored average fallback rate does not prove \(W\perp C\), does not control state-dependent selection after a cache miss, and does not establish the conditional path bounds.

\begin{theorem}[Harmless paths, revealing selector]\label{thm:selectorleak}
Let \(C\) be a uniform bit, choose \(W=C\), and let both path transcripts be constant. Then each path channel has zero maximal leakage conditional on its use, but the complete output reveals \(C\) through \(W\) and has one bit of maximal leakage.
\end{theorem}

Thus a path selector may dominate the privacy loss even when the path mechanisms themselves are perfect. When selection is secret- or state-dependent, the receipt must publish a direct selector/joint-channel bound or a valid composition supplied by the fallback-tax owner.\cite{series:fallbacktax}

\section{What an honest two-path receipt contains}\label{sec:receipt}
Each path line item binds:
\begin{itemize}[leftmargin=*]
\item secret, public side information, observer, complete projection, unit, horizon, and reset rule;
\item contacted-service or DHT surface, source/build/configuration digests, and mechanism knobs;
\item recognized channel-witness mode, witness digest, assumption evidence, and channel-level metric bound;
\item retry, timeout, failure, abort, stopping, and logging/extraction behavior;
\item a replay hook and drift invalidators.
\end{itemize}

The joint architecture additionally binds:
\begin{itemize}[leftmargin=*]
\item the selector implementation and whether \(W\) is visible;
\item evidence for secret independence, or a direct selector/joint bound;
\item correlation between selector, path channels, tiers, and adaptive state;
\item the combined budget and a fail-closed publication result.
\end{itemize}

A tier vector may organize contact, timing, and congestion projections, but every tier needs a channel witness and the joint release needs a factorization witness or direct joint bound.\cite{series:tieredobs}

\section{Operational priorities}\label{sec:priorities}
\begin{enumerate}[leftmargin=*]
\item Make the primary path's end-to-end reader channel explicit before optimizing its budget.
\item Treat fallback selection as a first-class output and test for secret/state dependence.
\item Cap fallback usage as availability and cost control, while proving rather than assuming its privacy effect.
\item Include downgrade, outage, timeout, and cache-miss behavior in hostile tests.
\item Recertify on contacted-service, source, configuration, selector, projection, or evidence drift.\cite{series:recert}
\end{enumerate}

\section{Takeaways}
\begin{itemize}[leftmargin=*]
\item Primary/fallback is an operational pattern, not a privacy theorem.
\item Each path needs a source-bound channel-level budget.
\item The selector can leak the entire secret even when both path outputs leak nothing.
\item A fallback rate or observation scalar is non-authorizing without its joint channel evidence.
\item The useful receipt publishes the boundary, assumptions, and blockers rather than merely the two path names.
\end{itemize}

\begin{thebibliography}{99}\small

\bibitem{series:bossfightB}
Anonymous.
\newblock AnonDHT in the Boss Fight Regime: Compiling a MaxL Budget into Lookup Knobs.
\newblock Boss Fight series Paper~B, 2026. (This archive.)


\bibitem{series:mceq}
Anonymous.
\newblock MC-EQ and CoverSketch: Destination Privacy via Markov-Chain Equalization in Anonymous DHT Lookups.
\newblock Release \& destination Paper~B, The-One-True-Archive (2026).

\bibitem{series:abomoinl}
Anonymous.
\newblock ABOM and OINL for Anonymity Claims.
\newblock Anonymity series Paper~C, The-One-True-Archive (2026).

\bibitem{series:release}
Anonymous.
\newblock Privacy as a Release Property: Proof-Carrying Anonymity Receipts for Anonymous DHTs.
\newblock Release \& destination Paper~A, The-One-True-Archive (2026).

\bibitem{series:recert}
Anonymous.
\newblock Disagreement-Gated Recertification.
\newblock Certified series Paper~C, The-One-True-Archive (2026).

\bibitem{series:threatwindows}
Anonymous.
\newblock Threat Models and Repetition Windows for Anonymous DHTs.
\newblock Synthesis~4, The-One-True-Archive (2026).

\bibitem{series:traceifaces}
Anonymous.
\newblock Trace Interfaces for Anonymous DHTs.
\newblock Synthesis~3, The-One-True-Archive (2026).


\bibitem{series:fallbacktax}
Anonymous.
\newblock Adaptive Fallback as a Witness Stream: The Path-Selection Tax and Honest Tiered Receipts.
\newblock Series synthesis note (Synthesis~14), The-One-True-Archive (2026).

\bibitem{series:obsatten}
Anonymous.
\newblock Observation Attenuation for Repeated-Use Privacy Budgets: Mixture/Subsampling Bounds and a Stub-Mix Interface.
\newblock Synthesis~11, The-One-True-Archive (2026).


\bibitem{series:tieredobs}
Anonymous.
\newblock Tiered Observation Vectors for Anonymous-DHT Receipts: Contact-Set, Timing, and Congestion Tiers with Conservative MaxL Composition.
\newblock Synthesis~15, The-One-True-Archive (2026).

\bibitem{series:state3}
Anonymous.
\newblock Closed-View Auditing and CPPC Manifests.
\newblock AnonDHT state series Paper~3, The-One-True-Archive (2026).

\bibitem{series:receiptitems}
Anonymous.
\newblock Receipt Line Items for Anonymity Claims: A Minimal Tuple Schema for Published Budgets.
\newblock Series synthesis note (Synthesis~12), The-One-True-Archive (2026).


\bibitem{series:deployedsurfaces}
Anonymous.
\newblock Deployed Observation Surfaces for Anonymous DHTs: Control-plane knobs in IPFS/libp2p delegated routing and OHTTP evolution.
\newblock Series synthesis note (Synthesis~30), 2026. (This archive.)

\end{thebibliography}

\end{document}
